%%%%%%%%%%%%%%%%%%%%%%%%%%%%%%%%%%%%%%%%%%%%%%%%%%%%%%%%%%
% FILE: environments.tex
%
% PURPOSE
% -------------------------------------------------------
% Defines reusable environments and macros used
% throughout the document.
%
% Includes two categories:
%
% 1. Float wrappers
%    Simplified commands for inserting figures,
%    tables and algorithms.
%
% 2. Documentation boxes
%    Visual blocks for definitions, notes,
%    warnings and examples.
%
% 3. Code blocks
%    Safe environment wrapper around lstlisting.
%
%%%%%%%%%%%%%%%%%%%%%%%%%%%%%%%%%%%%%%%%%%%%%%%%%%%%%%%%%%



% ==================================================
% FLOAT WRAPPERS
% ==================================================

% --------------------------------------------------
% FIGURE
% --------------------------------------------------

\newcommand{\figura}[4]{
\begin{figure}[H]
\centering
\includegraphics[width=\defaultfigurewidth]{#1}
\caption{#2}
\label{fig:#4}
\source{#3}
\end{figure}
}

% Usage:
% \figura{path}{caption}{source}{label}



% --------------------------------------------------
% TABLE
% --------------------------------------------------

\newcommand{\tabla}[4]{
\begin{table}[H]
\centering
\caption{#2}
\label{tab:#4}
#1
\source{#3}
\end{table}
}

% Usage:
% \tabla{table code}{caption}{source}{label}



% --------------------------------------------------
% ALGORITHM
% --------------------------------------------------

\newcommand{\algoritmo}[4]{
\begin{algorithm}[H]
\caption{#2}
\label{alg:#4}
#1
\source{#3}
\end{algorithm}
}

% Usage:
% \algoritmo{content}{caption}{source}{label}



% ==================================================
% DOCUMENTATION BOXES
% ==================================================
%
% Inspired by modern documentation systems such as
% Rust documentation, Python docs and MkDocs.
%

% --------------------------------------------------
% BASE STYLE
% --------------------------------------------------

\tcbset{
    docbox/.style={
        enhanced,
        colback=white,
        colframe=black,
        boxrule=0.6pt,
        arc=2pt,
        left=6pt,
        right=6pt,
        top=6pt,
        bottom=6pt,
        fonttitle=\bfseries,
        coltitle=black
    }
}



% --------------------------------------------------
% DEFINITION
% --------------------------------------------------

\newtcolorbox{definitionbox}[1][]{
    docbox,
    title={Definición: #1},
    colback=blue!5,
    colframe=blue!100!black
}



% --------------------------------------------------
% NOTE
% --------------------------------------------------

\newtcolorbox{notebox}{
    docbox,
    title={Nota},
    colback=gray!5,
    colframe=gray!60
}



% --------------------------------------------------
% WARNING
% --------------------------------------------------

\newtcolorbox{warningbox}{
    docbox,
    title={Advertencia},
    colback=red!3,
    colframe=red!70!black
}



% --------------------------------------------------
% EXAMPLE
% --------------------------------------------------

\newtcolorbox{examplebox}{
    docbox,
    title={Ejemplo},
    colback=green!3,
    colframe=green!60!black
}
